%% ============================================================
%% Section 7: Conclusions (V0.4)
%% Changes from V0.3.5: tone-down (R#10), limitations expanded (R#11),
%%   "most consequential" softened, §6 reference updated to 4 frontiers
%% ============================================================

\section{Conclusions}
\label{sec:conclusions}

This survey has traced the security and safety landscape of AI-Generated Content from its origins in content-level risks to the emerging frontier of autonomous agentic exploitation. A central meta-pattern emerges: \emph{temporal misalignment}. Attack surfaces evolve through a well-defined progression---direct prompt injection, jailbreaking, indirect prompt injection, and agentic exploitation via tool-invocation ecosystems---while defensive capabilities lag by months and governance frameworks by years. The first agentic AI products shipped in October 2024; the first government-level agentic governance framework appeared in January 2026, a 16-month gap during which AI agents with tool-invocation privileges were deployed at scale under no agent-specific standard~\cite{imda2026agentic}. The ``deploy first, retrofit security later'' pattern---traced concretely through the MCP deployment-to-governance timeline in Section~\ref{sec:threat-landscape} and documented throughout Sections~\ref{sec:threats}--\ref{sec:governance}---has extended from content risks to action risks with qualitatively higher consequences.

Key insights from each dimension reinforce this assessment. Generative capabilities have converged across modalities at near-zero marginal cost, with autonomous agency representing a qualitative shift from bounded output to unbounded action (Section~\ref{sec:capabilities}). The threat landscape spans content weaponization, model-level attacks, adversarial manipulation, and societal harms, with the attack surface evolution ladder (Table~\ref{tab:ladder}) from prompt injection to agentic exploitation constituting a highly consequential and currently under-governed vector (Section~\ref{sec:threats}). Technical countermeasures face fundamental limits---detection collapses under distribution shift, watermarking confronts information-theoretic impossibility results, alignment is superficial and fragile, and agentic defenses lack compositional guarantees (Section~\ref{sec:countermeasures}). Governance has fractured into three irreconcilable paradigms with zero coverage of agentic AI systems (Section~\ref{sec:governance}). Four prioritized research frontiers remain largely unsolved, with securing autonomous AI agents and unified multi-modal detection identified as the most urgent (Section~\ref{sec:future}).

\paragraph{Limitations.}
This survey has several limitations. \emph{Literature selection:} we surveyed peer-reviewed publications from major AI/security venues (2022--2026), supplemented by preprints from arXiv and institutional reports; no formal systematic review protocol (e.g., PRISMA) was applied, and inclusion decisions reflect the authors' judgment. \emph{Language and geographic bias:} our coverage is predominantly English-language and weighted toward Western and corporate perspectives on AI safety; important contributions in Chinese, Korean, and other languages may be underrepresented. \emph{Source reliability:} several key data points (e.g., industry-reported phishing statistics, MCP vulnerability audit results) originate from industry reports with undisclosed methodologies and potential conflicts of interest; the attack success rates cited throughout are drawn from individual studies and their reproducibility across settings has not been independently verified. \emph{Scope exclusions:} federated learning security, synthetic data training pipelines, and detailed red-teaming methodology are not covered; each merits dedicated survey treatment. \emph{Temporal constraints:} the knowledge cutoff of March 2026 constrains currency in a field evolving on monthly timescales, and the agentic AI security literature (predominantly 2025--2026) is too young for confident long-term trend assessment.

\paragraph{Looking forward.}
The asymmetry between offensive and defensive capabilities demands a paradigm shift from reactive to proactive security. Cross-disciplinary collaboration---spanning computer science, law, ethics, and public policy---is no longer aspirational but operationally necessary. The expansion from ``AI generates content'' to ``AI executes actions'' represents both a defining challenge and a critical research frontier for the coming decade.
